\chapter{Conclusion and Future Work} \label{chapter:ch7_conclusion}

\section{Conclusion}

This thesis asked whether collaborative-robot tasks that change frequently can be authored more directly by specifying spatial intent in the real workspace. Within prototype scope, the result is positive for selected task-oriented workflows. The developed system shows that poses and trajectories entered in the workspace can be combined with scene detection, task-specific interpretation, and shared robot execution to produce executable procedures on a real robot without exposing the operator to a full robot programming environment.

The main contribution is therefore a reusable task-oriented workflow rather than only a tracked input device or two isolated demonstrations. The thesis defines a use-case boundary in which selected tasks can be authored, confirmed into executable commands, executed, visualized, and persisted through one end-to-end workflow. This workflow is supported by a shared modular runtime with explicit interfaces between sensing, authoring, interpretation, visualization, and execution, which allows different use cases to reuse the same infrastructure. The result is demonstrated on two representative task families: pick-and-place and seam welding.

At the same time, the result remains a proof of concept. The evaluation is qualitative, physical validation is limited, and the thesis does not include a formal user study, a baseline comparison against established programming workflows, or a quantitative characterization of tracking accuracy, repeatability, or productivity. Within that boundary, seam welding is the clearest practical fit for the current prototype because the authored spatial input maps naturally to a trajectory-centric task. Pick-and-place is more important as flexibility evidence: it shows that authored intent can be combined with refreshed scene information at execution time, but it depends more strongly on the narrow scene model used in the prototype.

\section{Near-Term Future Work}

The first near-term step is to complete the sensing side of the prototype. This includes more robust pen tracking, better calibration diagnostics, and fuller use of the sensing already present on the pen side, especially IMU data. The scene pipeline should also move beyond the current fiducial-box representation toward a richer and better characterized scene estimate, with stronger handling of tracking stability, scene uncertainty, and repeated scene acquisition in the real workspace.

The second near-term step is to complete the execution layer. A more mature version of the system would need a clearer abstraction for process tools, explicit motion validation, and collision checking for authored tasks. Simulated execution would also be valuable, both for safer inspection of authored procedures and for development of use cases before real-robot execution. Where a task does not require runtime scene adaptation, it would also be useful to explore a path from authored commands to a lower-level export or compilation target on the robot side.

The third near-term step is stronger validation of the current prototype. The present work would benefit from broader physical testing and from bounded characterization of repeatability, tracking accuracy, and scene-estimation stability in the demonstrated setup. Cross-platform completion and some control-plane cleanup also belong in this group. These are not new research directions by themselves, but they would make the current system a more reliable platform for continued evaluation.

\section{Longer-Term Research Directions}

Beyond prototype completion, one natural direction is richer input for spatial authoring. The architecture is not tied to the current pen-based interface, so alternative inputs such as hand tracking, VR or AR controllers, or combinations of speech and gesture could be explored within the same task-oriented model. This would help clarify which kinds of input are most suitable for different use cases and how much spatial intent can be specified naturally without losing control over the resulting robot behavior.

Another direction is to study how this authoring model fits existing robot-programming practice. One possibility is deeper integration with teach-pendant or controller workflows, so that spatial authoring becomes one programming tool among others rather than a separate frontend. Another is a formal user study focused on when this approach is actually useful: for which task classes it reduces setup friction, where its limits remain too restrictive, and how operators understand the boundary between prepared use cases and more general robot programming.

Finally, the prototype should be exercised on a broader range of task families and operating conditions. The current thesis suggests that use-case-oriented spatial authoring is most promising where the physical task is repetitive enough for robot execution to be worthwhile, but the concrete task definition changes often enough that repeated low-level programming becomes burdensome. Extending the work toward richer welding scenarios, additional fabrication or surface-processing tasks, and more varied human-in-the-loop workflows would provide a stronger basis for judging that applicability boundary. The present prototype provides a bounded starting point for that work, but broader claims depend on the technical completion and evaluation outlined above.
